\documentclass{article}
\usepackage{graphicx} % Required for inserting images

\title{Optimal Path for Package distribution}
\author{Romano Brentani}

\date{April 2026}

\begin{document}

\maketitle

\section{Introduction}

\section{Domain}
We consider a warehouse where a robot must pick up $N$ packages and deliver them to different locations while navigating around obstacles. The warehouse is represented as a grid and the problem is to figure out in which order the packages should be processed and stored. 
Each package must be stored in one storage and the Agent can store one package at the time. The goal of the Agent is to determine how to navigate efficiently across all deliveries without running into a obstacle. We model this problem as a Markov Decision Process (MDP), where the agent learns through trial and error to find the best delivery strategy.

\section{Methodology}

\section{Experiment Design}

\subsection{Objective}

The objective of the experiments is to evaluate whether reinforcement learning methods enable a warehouse robot to efficiently distribute and store incoming packages in a constrained environment. The agent must optimize navigation and task sequencing while operating in the presence of obstacles and environmental uncertainty.

Four reinforcement Learning approaches will be compared:

\begin{itemize}
    \item Q-Learning,
    \item SARSA,
    \item Value Iteration,
    \item Policy Iteration
    \item Deep Q-Learning
\end{itemize}

The experiments aim to answer the following research questions:

\begin{enumerate}
    \item Does reinforcement learning methods learning improve warehouse efficiency compared to heuristic solutions?
    \item Which algorithm achieves the best performance under varying environmental conditions?
    \item How robust are learned policies across different warehouse configurations?
\end{enumerate}

\subsection{Environment Setup}

The warehouse is modeled as a discrete grid-environment representing storage locations, free navigation space, and obstacles.

In each experiment-setup, the robot receives $N$ packages that must be delivered to designated storage positions. The number $N$ can be configured individual by the user. 

The environment state is defined as:
\begin{itemize}
    \item robot position,
    \item remaining packages,
    \item storage occupancy,
    \item obstacle configuration.
\end{itemize}

The agent can execute the following actions:

\begin{itemize}
    \movement actions (up, down, left, right),
    \package pickup,
    \package storage.
\end{itemize}

Obstacles can be static or dynamically added during execution to simulate realistic warehouse disruptions such as blocked paths or fallen shelves.

\subsection{Reward Design}

The goal is to complete the task described above efficient.
In order to achieve that a reward function is defined. The reward function is designed to encourage efficient task completion:

\begin{equation}
R = -c_{m} - c_{t} + r_{s}
\end{equation}

where:
\begin{itemize}
    \item $c_{m}$ represents unnecessary movement,
    \item $c_{t}$ penalizes longer completion times,
    \item $r_{s}$ rewards the successful package storage.
    \item $c_{col}$ is the punishment for collide with a obstacle
\end{itemize}

A successful learning method aims to minimize the total completion time, penalty factors and tries to maximize the cumulative reward.

\subsection{Baseline Methods}

To assess whether reinforcement learning provides measurable benefits, we compare against established non-learning baselines:

\begin{itemize}
    \item \textbf{Shortest-path heuristic:} path planning using deterministic search (e.g., A* or Dijkstra’s Algorithm) combined with fixed package ordering,
    \item \textbf{Greedy strategy:} always storing the nearest available package,
    \item \textbf{Random policy:} lower-bound performance reference.
\end{itemize}

These baselines are commonly used rule-based warehouse strategies.

\subsection{Compared Algorithms}

\paragraph{Q-Learning}
Q-Learning is a simulation-based value iteration algorithm. This algorithm uses a trivial empirical MDP model. That means the true MDP model is replaced by an estimate MDP model. 
Q-Learning is a temporal-difference method that tries to learns optimal action-values independent of the policy behavior. It is expected to achieve strong performance in stationary environments.

\paragraph{SARSA}
An on-policy temporal-difference algorithm that updates value estimates using the executed policy. Due to its conservative updates, it is expected to perform robustly in stochastic environments.

\paragraph{Value Iteration}
Value iteration is a simple extension of
the backwards induction algorithm to the infinite horizon case. 
Mathematically, Value iteration corresponds to repeated application of the Bellman operator. 

\paragraph{Policy Iteration}
policy iteration attempts to iteratively improve a given
policy, rather than a value function. Starting with an arbitrary policy, at each
iteration it first computes the value of the current policy. In a second step called policy improvement the policy is updated by choosing the policy that is greedy with respect to the value function computed in the evaluation step.

\subsection{Evaluation Metrics}

Performance is evaluated using the following metrics:

\textbf{Efficiency Metrics}

\begin{itemize}
    \item total task completion time,
    \item total travel distance,
    \item number of executed actions,
    \item energy consumption proxy.
\end{itemize}

\textbf{Learning Metrics}
\begin{itemize}
    \item cumulative reward per episode,
    \item convergence speed,
    \item policy stability after convergence.
\end{itemize}

\textbf{Operational Metrics}
\begin{itemize}
    \item collision rate,
    \item failed storage attempts,
    \item average decision latency.
\end{itemize}

An overall efficiency score is defined as:

\begin{equation}
E= \frac{{Packages Stored}}{{TimeSteps}}
\end{equation}

\subsection{Experimental Procedure}

Each algorithm is trained over multiple episodes until convergence.

The evaluation procedure consists of:

\begin{enumerate}
    \item training policies in a fixed warehouse configuration,
    \item testing on unseen environment layouts,
    \item repeating experiments across multiple random seeds,
    \item reporting averaged performance statistics.
\end{enumerate}

Training and evaluation environments are separated to measure generalization capability.

\subsection{Expected Contexts of Improvement}

Reinforcement learning approaches are expected to outperform heuristic baselines in scenarios where:

\begin{itemize}
    \item obstacle layouts create complex navigation constraints,
    \item package ordering significantly affects efficiency,
    \item environmental conditions change dynamically,
    \item global optimization is required beyond local greedy decisions.
\end{itemize}

\subsection{Robustness Evaluation}

Robustness is evaluated by systematically varying the problem instance:

\begin{itemize}
    \item different numbers of packages $N$,
    \item varying warehouse sizes,
    \item randomized obstacle placements,
    \item dynamically appearing obstacles,
    \item stochastic action outcomes.
\end{itemize}

Robustness is measured through performance variance and degradation under increasing task complexity.

\end{document}
